\section{INTRODUCCIÓN}

El crecimiento acelerado de las plataformas de comercio electrónico y los portales de gestión pública en Bolivia ha configurado un entorno digital caracterizado por una alta densidad de datos. No obstante, a diferencia de otros ecosistemas internacionales consolidados, el ámbito digital boliviano presenta una marcada asimetría informativa y una notable opacidad en la publicación de datos abiertos estructurados. Las principales cadenas de retail minorista y cadenas farmacéuticas del país carecen de interfaces de programación de aplicaciones (API) públicas y accesibles para la ciudadanía o la comunidad académica. Esta limitación restringe la capacidad de monitorear oportunamente la fluctuación de precios de la canasta familiar y evaluar la dispersión de costos de medicamentos esenciales en los centros urbanos del eje central: La Paz, Cochabamba y Santa Cruz.

De manera simultánea, las plataformas gubernamentales como el Sistema de Contrataciones Estatales (SICOES) y el Observatorio Agroambiental y Productivo (OAP) concentran volúmenes sustanciales de convocatorias públicas y series de abastecimiento agropecuario. Sin embargo, su consumo automatizado se ve obstaculizado por esquemas de búsqueda tradicionales basados exclusivamente en concordancia exacta de cadenas de texto y arquitecturas web protegidas por cortafuegos de aplicaciones web (WAF, por sus siglas en inglés).

Para responder a estos desafíos tecnológicos y sociales, el presente proyecto académico desarrolla e implementa un \textbf{Sistema Integral de Minería Web Multidimensional}, articulando de manera armónica las tres ramas clásicas de la disciplina conceptualizadas por \citet{liu2011web} y \citet{han2012data}:

\begin{enumerate}[leftmargin=*, label=\textbf{\arabic*.}]
    \item \textbf{Minería de Contenido Web (\textit{Web Content Mining}):} Construcción y despliegue de agentes de extracción automatizada y resiliente basados en Node.js y Playwright, diseñados para superar barreras dinámicas de representación del lado del cliente (\textit{Client-Side Rendering}), desplazamiento infinito (\textit{infinite scrolling}) y mecanismos de protección perimetral basados en desafíos biométricos y de comportamiento humano (WAF Radware Bot Manager y hCAPTCHA). Este subsistema opera sobre los catálogos en línea de Hipermaxi, Farmacorp, Farmacias Chávez, el portal de licitaciones de SICOES y el observatorio de precios del OAP.
    
    \item \textbf{Minería de Estructura Web (\textit{Web Structure Mining}):} Extracción, modelado y análisis de la topología de red de contrataciones estatales en el departamento de La Paz. El sistema construye grafos bipartitos dirigidos conectando entidades públicas convocantes con rubros comerciales y proveedores adjudicados, aplicando formalmente los algoritmos \textbf{PageRank} \citep{page1998anatomy} y \textbf{HITS} (\textit{Hyperlink-Induced Topic Search}) \citep{kleinberg1999authoritative} para identificar concentraciones de poder licitador y autoridades temáticas en compras públicas.
    
    \item \textbf{Minería de Uso Web (\textit{Web Usage Mining}):} Instrumentación de telemetría de eventos, tiempos de respuesta y patrones de interacción humana mediante el asistente conversacional oficial desplegado en Telegram (\texttt{@ChatDM7200bot}) \citep{srivastava2000web, cooley1999data}. Se capturan trazas operacionales para optimizar la navegación faceted de los usuarios y enriquecer los flujos de consulta en lenguaje natural.
\end{enumerate}

El ecosistema integra una arquitectura de datos en la nube sobre \textbf{Neon Serverless PostgreSQL}, estructurada mediante un modelo estrella relacional que combina un catálogo maestro inmutable de productos (\textit{Master Catalog}), series de tiempo particionadas de forma declarativa por rangos mensuales (\textit{Time-Series Partitioning}) y triggers reactivos PL/pgSQL que calculan en tiempo de inserción las variaciones de precio ($\Delta\text{Bs}$ y $\Delta\%$). Adicionalmente, el pipeline multimedia desacopla los activos visuales efímeros de las plataformas origen mediante un almacenamiento persistente en \textbf{Neon Object Storage S3}, sirviendo los activos a través de una plataforma web interactiva construida en \textbf{Next.js 14} y complementada con búsqueda semántica vectorial densa impulsada por modelos de inteligencia artificial \textbf{NVIDIA NIM} (1024 dimensiones) y cartografía geoespacial con \textbf{OpenStreetMap}.

El presente informe detalla el marco teórico, metodológico y práctico del proyecto, exponiendo el diseño arquitectónico, las evidencias visuales y las métricas de rendimiento auditadas en el entorno operativo.
